\chapter{Problem 4 -- Arsitektur Microservice dan Flexible Storage}

\begin{ringkas}
Sistem dipecah menjadi \emph{container} yang masing-masing hanya punya satu alasan untuk berubah: \emph{mock} sumber, Aggregator,
Auth Service, Client-Facing API, dan \emph{consumer}. Canonical Store memakai PostgreSQL dengan kolom JSONB, sehingga
\emph{field} baru seperti \code{confidence_level} tidak membutuhkan perubahan tabel. Jaringan Docker \emph{internal} membuat
database secara fisik tidak terjangkau oleh \emph{service} selain Aggregator.
\end{ringkas}

\begin{catatan}[Status bukti Problem 4]
Ketiga kriteria dibuktikan oleh satu skrip, \cmd{make demo-p4}. Pada langkah A, \code{pvmbg-mock} dihentikan, dibangun
ulang, dan dijalankan lagi sementara \emph{container} lain tetap berjalan. Pada langkah B sampai D, \emph{field} baru masuk ke
PostgreSQL JSONB tanpa migrasi dan data lama tetap terbaca. Pada langkah E, terbukti hanya Aggregator yang dapat
menjangkau Canonical Store.
\end{catatan}

\section{Rancangan solusi dan mekanisme}

\subsection{Pembagian service}

Batas antarservice ditarik dengan satu aturan: satu \emph{container} punya satu alasan untuk berubah, dan setiap data
punya satu pemilik. Perbaikan untuk satu alasan tidak boleh memaksa \emph{container} lain di-\emph{rebuild}.

\begin{table}[H]
\centering\small
\begin{tabularx}{\textwidth}{@{}L{3.8cm}L{4.4cm}Y@{}}
\toprule
\textbf{\emph{Container}} & \textbf{Tanggung jawab tunggal} & \textbf{Berubah ketika}\\
\midrule
\code{bmkg-mock}, \code{pvmbg-mock} & Menyimulasikan lembaga lain dan bukan bagian BNPB & kontrak sumber berubah (di luar kendali BNPB)\\
\code{aggregator} & Mengambil data, memetakannya, memiliki Canonical Store, dan mempublikasikan \emph{event} & aturan pemetaan atau \emph{ingest} berubah\\
\code{auth-service} & Menerbitkan dan memperbarui \emph{token} klien & kebijakan autentikasi berubah\\
\code{client-api} & Menyajikan data serta menegakkan hak akses dan \emph{rate limit}, tanpa \emph{state} & kebijakan \emph{field} atau \emph{rate limit} berubah\\
\code{consumer-*} & Melayani satu kebutuhan pihak yang berlangganan (\emph{dashboard}, notifikasi) & kebutuhan pihak tersebut berubah\\
\code{postgres-canonical}, \code{kafka} & Infrastruktur penyimpanan dan distribusi \emph{event} & versi atau konfigurasi infrastruktur berubah\\
\bottomrule
\end{tabularx}
\caption{Batas \emph{service} dan alasan masing-masing berubah.}
\end{table}

Spesifikasi memberi contoh \emph{perbaikan kecil pada autentikasi}. Perbaikan seperti itu hanya menyentuh
\code{auth-service} (atau \code{client-api} untuk bagian pemeriksaannya). Aggregator, \emph{mock}, dan \emph{consumer} tidak ikut
di-\emph{rebuild}.

\subsection{Aturan kepemilikan penyimpanan}

Setiap penyimpanan hanya dimiliki oleh satu \emph{service}. Komponen lain meminta data lewat API milik pemiliknya.

\begin{table}[H]
\centering\small
\begin{tabularx}{\textwidth}{@{}L{3.6cm}L{3.0cm}Y@{}}
\toprule
\textbf{Penyimpanan} & \textbf{Pemilik} & \textbf{Cara komponen lain memperoleh data}\\
\midrule
Canonical Store (PostgreSQL) & \code{aggregator} & API internal Aggregator (\code{/internal/*}) dengan \code{X-Service-Key}\\
\emph{Hash} \emph{refresh} \emph{token} (SQLite) & \code{auth-service} & \emph{endpoint} \code{/v1/token} dan \code{/v1/refresh}\\
\emph{Topic} Kafka \code{hazard.events} & \code{aggregator} sebagai publisher & men-\emph{subscribe} sebagai \emph{consumer group}\\
Kumpulan \code{event_id} di \emph{consumer} & masing-masing \emph{consumer} & tidak dibagikan\\
\bottomrule
\end{tabularx}
\caption{Kepemilikan penyimpanan.}
\end{table}

\subsection{Penyimpanan yang fleksibel: PostgreSQL dengan JSONB}

Tabel \code{hazards} hanya memakai kolom tetap untuk identitas dan \emph{field} kanonik. Semua yang bisa berubah masuk ke satu
kolom \code{attributes} bertipe JSONB.

\begin{lstlisting}[language=SQL]
CREATE TABLE IF NOT EXISTS hazards (
  hazard_id     text PRIMARY KEY,
  source        text NOT NULL,     source_ref_id text NOT NULL,
  hazard_type   text NOT NULL,     severity      text NOT NULL,
  area_name     text NOT NULL,
  latitude      double precision,  longitude     double precision,
  occurred_at   timestamptz NOT NULL,  ingested_at timestamptz NOT NULL,
  updated_at    timestamptz NOT NULL,  version     integer NOT NULL,
  attributes    jsonb NOT NULL DEFAULT '{}'      -- field dinamis: tanpa DDL saat field baru muncul
);
\end{lstlisting}

Potongan di atas sudah disederhanakan. Definisi aslinya (\code{adapters/postgres/repo.go}) juga memuat indeks dan aturan
pengurutan teks untuk \code{hazard_id}, tetapi tidak ada kolom lain. Perintah \cmd{CREATE ... IF NOT EXISTS} dijalankan
\emph{sekali saat container start}, bukan sebagai langkah migrasi setiap kali sumber berubah. Repositori tidak memuat
direktori migrasi maupun berkas SQL terpisah. Ketika PVMBG mulai mengirim \code{confidence_level}, Aggregator tidak
di-\emph{restart} dan tidak ada DDL yang dijalankan, karena \emph{field} itu langsung masuk ke \code{attributes}.

\subsection{Isolasi jaringan sebagai bukti kepemilikan}

Klaim bahwa hanya Aggregator yang memegang database tidak dibiarkan sebagai janji. Compose membuat empat jaringan, dan
tiga di antaranya bersifat \emph{internal} (tidak bisa diakses dari luar dan tidak bisa dipakai untuk membuka port ke
host).

\begin{table}[H]
\centering\small
\begin{tabularx}{\textwidth}{@{}L{3.3cm}L{2.2cm}Y@{}}
\toprule
\textbf{Jaringan} & \textbf{Internal?} & \textbf{Anggota}\\
\midrule
\code{app-net} & tidak & \emph{mock}, aggregator, auth-\emph{service}, client-api, \emph{consumer} (hanya port \code{/health} dan \code{/stats})\\
\code{canonical-net} & ya & hanya aggregator dan postgres-canonical\\
\code{broker-net} & ya & kafka, aggregator (publisher), \emph{consumer}\\
\code{auth-net} & ya & auth-\emph{service}\\
\bottomrule
\end{tabularx}
\caption{Jaringan Docker dan anggotanya.}
\end{table}

\section{Alasan keputusan dan trade-off}

\subsection{Batas service: mengapa dipecah seperti ini}\label{sec:p4-batas}

Sistem ini bisa dipecah menjadi \emph{service} dengan beberapa cara. Pembagian yang kami pilih dibandingkan dengan alternatifnya di bawah.

\begin{table}[H]
\centering\small
\begin{tabularx}{\textwidth}{@{}L{3.9cm}Y@{}}
\toprule
\textbf{Pembagian alternatif} & \textbf{Pertimbangan}\\
\midrule
Monolit modular (satu \emph{container}) & Paling sederhana, tetapi inilah kondisi yang menjadi masalah pada spesifikasi: perbaikan autentikasi memaksa seluruh sistem di-\emph{rebuild}.\\
Dipecah per lapisan teknis (poller, \emph{mapper}, penyimpan, publisher) & Terlalu halus. Hazard dan baris \emph{outbox} harus ditulis dalam \emph{satu transaksi}, dan jika dipisah, transaksi itu menjadi transaksi terdistribusi. Pemetaan juga membutuhkan hasil \emph{polling} yang sama, sehingga memisahkannya hanya menambah satu lompatan jaringan tanpa manfaat \emph{deploy} mandiri.\\
Dipecah per sumber (adapter BMKG dan adapter PVMBG masing-masing di \emph{container} sendiri) & Menarik, karena perubahan skema PVMBG hanya me-\emph{rebuild} adapter PVMBG. Tidak diambil pada M1 sebab setiap adapter butuh kontrak, \emph{container}, dan jalur ke \emph{store} tersendiri, sedangkan di dalam Aggregator kedua sumber sudah berjalan di putaran terpisah. Kandidat untuk M2.\\
\textbf{Dipecah per alasan berubah (dipilih)} & Sumber, Aggregator (pemilik data), autentikasi, penyajian, dan \emph{consumer} dapat di-\emph{rebuild} secara mandiri. Pasangan yang harus konsisten secara transaksi tetap berada di satu \emph{container}.\\
\bottomrule
\end{tabularx}
\caption{Alternatif pembagian \emph{service}.}
\end{table}

Konsekuensinya, mengubah aturan pemetaan berarti me-\emph{rebuild} Aggregator, termasuk \emph{polling}-nya. Hal ini tidak berbahaya
karena penanda kemajuan \emph{polling} (\emph{watermark}) tersimpan di database dan Aggregator tidak bergantung pada urutan
\emph{start} (tanpa \code{depends_on} ke \emph{mock} maupun \emph{store}, dan koneksi dicoba ulang sampai 90~detik). Auth Service,
Client-Facing API, dan \emph{consumer} tidak ikut tersentuh.

\subsection{Pilihan penyimpanan}\label{sec:p4-storage}

Kebutuhan utamanya adalah \emph{field} baru dari sumber tidak boleh memaksa migrasi tabel. Tabel berikut membandingkan pilihan
yang masuk akal.

\begin{table}[H]
\centering\small
\begin{tabularx}{\textwidth}{@{}L{3.1cm}YY@{}}
\toprule
\textbf{Alternatif} & \textbf{Kelebihan} & \textbf{Konsekuensi bagi sistem ini}\\
\midrule
Relasional dengan kolom tetap & Skema tegas dan \emph{query} kuat & Setiap \emph{field} baru membutuhkan migrasi dan \emph{deploy} bersama, sehingga masalah Problem~1 muncul kembali lewat jalur \emph{storage}. Ditolak.\\
Tabel atribut terpisah (EAV) & Fleksibel tanpa DDL & \emph{Query} menjadi rumit dan lambat (banyak \emph{join}) dan tipe data hilang. Lebih buruk daripada JSONB.\\
Document \emph{store} (misalnya MongoDB) & Cocok untuk dokumen tanpa skema & Menambah satu teknologi baru. Menulis hazard dan \emph{outbox} secara atomik membutuhkan transaksi lintas dokumen (perlu \emph{replica set}), yang lebih berat untuk satu \emph{instance} di Compose.\\
\emph{Key-value} (misalnya Redis) & Sangat cepat & Tidak mendukung \emph{query} berurutan dan paginasi berbasis kunci yang kami butuhkan, serta tidak ada transaksi hazard dan \emph{outbox}.\\
\textbf{PostgreSQL dengan JSONB (dipilih)} & Kolom tetap untuk data yang pasti dan JSONB untuk yang berubah, transaksi ACID, indeks dan \emph{query} lengkap, serta satu teknologi untuk hazard dan \emph{outbox} & Lihat konsekuensi di bawah.\\
\bottomrule
\end{tabularx}
\caption{Perbandingan alternatif penyimpanan.}
\end{table}

Konsekuensi dari pilihan ini:
\begin{itemize}
 \item Keuntungan: \emph{field} baru tidak membutuhkan migrasi, hazard dan entri \emph{outbox} ditulis secara atomik (dasar bagi Problem~5), dan \emph{keyset pagination} didukung oleh indeks \cmd{(hazard_type, occurred_at, hazard_id)}.
 \item Kerugian: isi \code{attributes} tidak divalidasi tipenya dan belum diindeks, sehingga API belum mendukung penyaringan berdasarkan nilai atribut. Jika kelak dibutuhkan, indeks GIN dapat ditambahkan tanpa mengubah kontrak.
 \item Keputusan kecil yang disengaja: kolom \code{outbox.payload} bertipe \code{json}, bukan \code{jsonb}, agar pesan tersimpan persis seperti saat dibuat (JSONB menata ulang urutan kunci).
\end{itemize}

\section{Batasan yang disadari}

\begin{itemize}
 \item Demo \emph{rebuild} memakai \code{pvmbg-mock}, bukan Aggregator. \emph{Mock} itu dipilih karena paling aman dimatikan saat demo berjalan. Me-\emph{rebuild} Aggregator berarti \emph{polling} berhenti sebentar, dan hal itu tidak diperagakan oleh skrip ini.
 \item Selama PVMBG di-\emph{rebuild}, data vulkanik tidak bertambah sampai \emph{polling} pulih. Data seismik dan \emph{record} lama tetap tersedia dari Canonical Store.
 \item Aggregator masih menggabungkan beberapa tugas (\emph{polling}, pemetaan, kepemilikan \emph{store}, \emph{relay}, dan API internal). Ini keputusan yang disengaja (alasannya di atas), tetapi berarti \emph{container} tersebut relatif besar.
 \item \code{pkg/obs} dipakai bersama oleh semua \emph{service} (\emph{log}, \emph{correlation ID}, \emph{health}). Isinya hanya utilitas teknis tanpa logika bisnis, sehingga sesuai aturan pada spesifikasi. Meski begitu, perubahan di sana tetap berarti semua \emph{service} perlu di-\emph{rebuild} agar memakainya.
 \item PostgreSQL dan Kafka masing-masing hanya satu \emph{instance} tanpa replika. Ketahanan terhadap kegagalan bersamaan menjadi bahan untuk M2.
 \item Port API internal Aggregator dibuka ke \emph{loopback} host (127.0.0.1:8083) untuk operator dan skrip demo, dan tidak dipakai oleh klien.
 \item Karantina disimpan di memori, bukan di Canonical Store.
\end{itemize}

\section{Bukti penyelesaian}

Setiap kriteria pada spesifikasi dan bukti yang menjawabnya dirangkum di bawah.

\begin{table}[H]
\centering\small
\begin{tabularx}{\textwidth}{@{}L{4.3cm}YL{1.9cm}@{}}
\toprule
\textbf{Kriteria pada spesifikasi} & \textbf{Bukti} & \textbf{Status}\\
\midrule
Setiap komponen berjalan sebagai \emph{container} terpisah & Setiap \emph{service} punya \code{Dockerfile} sendiri dan satu perintah (\cmd{docker compose up --build}) menyalakan semuanya. Daftar \emph{container} ada di Tabel~\ref{tab:container}. & \lulus\\
Satu \emph{container} dapat dihentikan, di-\emph{rebuild}, dan dijalankan ulang tanpa me-\emph{restart} \emph{container} lain, dan sistem tetap melayani permintaan lain & \cmd{make demo-p4}, langkah A: \code{pvmbg-mock} dihentikan, di-\emph{rebuild} (\cmd{docker compose build}), lalu dijalankan lagi dengan \cmd{docker compose up -d --no-deps}. Selama proses itu, permintaan khusus data BMKG lewat Client-Facing API tetap berstatus 200. ID \emph{container} dan waktu \emph{start} delapan \emph{container} lain dicatat sebelum dan sesudah, dan harus identik. & \lulus\\
HazardEvent tersimpan tanpa migrasi saat \emph{field} baru muncul, dengan \emph{record} lama dan baru berdampingan & \cmd{make demo-p4}, langkah B sampai D. Laporan lama (tanpa \code{confidence_level}) disimpan, skema PVMBG diubah saat sistem berjalan, lalu laporan baru disimpan. Keduanya ada di PostgreSQL dan terbaca lewat Aggregator, dan atribut laporan lama tidak berubah. \emph{Fingerprint} kolom tabel \code{hazards} sama sebelum dan sesudah, begitu pula waktu \emph{start} Aggregator, sehingga tidak ada DDL dan tidak ada \emph{restart}. Demo Problem~1 menunjukkan hal serupa pada skala lebih besar (5 hazard dengan dan 28 tanpa \emph{field} itu). & \lulus\\
Canonical Store hanya dapat diakses langsung oleh Aggregator & \cmd{make demo-isolation}, yang juga dijalankan sebagai langkah E pada \cmd{make demo-p4}. Cuplikannya ada di bawah. & \lulus\\
\bottomrule
\end{tabularx}
\caption{Kriteria Problem~4 dan buktinya.}
\end{table}

\subsection*{Isolasi Canonical Store (\cmd{make demo-isolation})}
\begin{lstlisting}
== [1] Who is attached to which network (docker network inspect) ==
   bnpb_app-net: aggregator auth-service bmkg-mock client-api consumer-dashboard consumer-notifier pvmbg-mock
   bnpb_canonical-net: aggregator postgres-canonical
   bnpb_broker-net: aggregator consumer-dashboard consumer-notifier kafka

== [2] PostgreSQL (Canonical Store) is reachable ONLY from the Aggregator's network ==
   networks the database container is attached to: bnpb_canonical-net
   PASS canonical-net contains exactly: aggregator + postgres-canonical
   PASS postgres publishes no port to the host
   PASS container on app-net (where mocks / client-facing API live): postgres-canonical:5432 is unreachable
   PASS container on broker-net (consumers): postgres-canonical:5432 is unreachable
   PASS control: a container on canonical-net (i.e. the Aggregator's network) does reach it

== [3] Kafka is reachable only from broker-net (producer + consumers), not from app-net ==
   PASS container on app-net: kafka:9092 is unreachable
   PASS control: a container on broker-net does reach kafka:9092

== [4] Other components get hazards only through the Aggregator API ==
   PASS the Aggregator API demands its service key (401 without, 403 with a wrong one)

ISOLATION PASS
\end{lstlisting}

Baris \emph{control} penting karena membuktikan bahwa penolakan di atas \emph{bukan} akibat pengujian yang rusak.
\emph{Container} percobaan di jaringan yang sama dengan Aggregator berhasil terhubung, sedangkan \emph{container} di jaringan lain
tidak.

\subsection*{Field baru tanpa migrasi (\cmd{make demo-p4}, langkah B--D)}
\begin{lstlisting}
== [2] B. Store an old-schema volcanic report ==
   old row attributes: {"eruption_count_24h": 3, "ash_column_height_m": 2730}
   hazards table column fingerprint before: 4555805e1c5790a6fff3c9ce2923df04
   PASS pvmbg-VR-0c2b58-0025 stored without confidence_level in PostgreSQL JSONB

== [3] C. Change PVMBG's schema live and store a new report without a migration ==
   new row attributes: {"confidence_level": 0.82, "eruption_count_24h": 15, "ash_column_height_m": 1822.6}
   hazards table column fingerprint after: 4555805e1c5790a6fff3c9ce2923df04
   PASS pvmbg-VR-0c2b58-0028 contains confidence_level; table columns and Aggregator start time did not change

== [4] D. Read old and new records side by side through the Aggregator ==
   PASS old=VR-0c2b58-0025 and new=VR-0c2b58-0028 coexist and are readable; old attributes are unchanged
\end{lstlisting}

\emph{Fingerprint} yang sama sebelum dan sesudah berarti daftar kolom tabel tidak berubah. \emph{Field} baru hanya muncul di
dalam kolom \code{attributes}.

\subsection*{Me-rebuild satu service (\cmd{make demo-p4}, langkah A)}
\begin{lstlisting}
== [1] A. Rebuild only PVMBG while unrelated services keep serving requests ==
   unrelated containers before rebuild:
   aggregator|8a9714f802b741b729024cd242a70a9332a1900012139de893d81bf779c3ad87|2026-10-08T09:58:31.792153279Z
   auth-service|af336456bb62348ccbfd24fc36414365af32afedd99cbcc8d88446ba54e27120|2026-10-08T09:58:31.736886278Z
   bmkg-mock|2f0e674affad290bbab61f78f1e40eddaf9fbdcd283cf44f945f761f19247ca5|2026-10-08T09:58:31.67503029Z
   client-api|92d5f361c8065ff9627d59efdeff89666a20d68523f0424cb21d886bc80ebb49|2026-10-08T09:58:37.840711994Z
   consumer-dashboard|7668604c282376fb3532f09849c904c6a46f26b6a73474a9993c9f052a8a4643|2026-10-08T09:58:39.594083631Z
   consumer-notifier|1aad8bc2ce8821326644d5ffbc9bd4702a5d5380bd384775992ae2b048accf61|2026-10-08T10:02:26.004276246Z
   kafka|fafda708f27fdd3b3905998f5f39ed4764074d3597038709bc84546cc439ac57|2026-10-08T09:58:31.663436421Z
   postgres-canonical|20fad02bbd16a6feedbfcc2b1edcbed8b1f5cffd88676cf812ca9b17861e608d|2026-10-08T09:58:31.609251372Z
   unrelated containers after rebuild:
   aggregator|8a9714f802b741b729024cd242a70a9332a1900012139de893d81bf779c3ad87|2026-10-08T09:58:31.792153279Z
   auth-service|af336456bb62348ccbfd24fc36414365af32afedd99cbcc8d88446ba54e27120|2026-10-08T09:58:31.736886278Z
   bmkg-mock|2f0e674affad290bbab61f78f1e40eddaf9fbdcd283cf44f945f761f19247ca5|2026-10-08T09:58:31.67503029Z
   client-api|92d5f361c8065ff9627d59efdeff89666a20d68523f0424cb21d886bc80ebb49|2026-10-08T09:58:37.840711994Z
   consumer-dashboard|7668604c282376fb3532f09849c904c6a46f26b6a73474a9993c9f052a8a4643|2026-10-08T09:58:39.594083631Z
   consumer-notifier|1aad8bc2ce8821326644d5ffbc9bd4702a5d5380bd384775992ae2b048accf61|2026-10-08T10:02:26.004276246Z
   kafka|fafda708f27fdd3b3905998f5f39ed4764074d3597038709bc84546cc439ac57|2026-10-08T09:58:31.663436421Z
   postgres-canonical|20fad02bbd16a6feedbfcc2b1edcbed8b1f5cffd88676cf812ca9b17861e608d|2026-10-08T09:58:31.609251372Z
   PASS BMKG-only requests succeeded; IDs and start times of all other containers stayed fixed
\end{lstlisting}

\section{Jawaban}

\begin{jawaban}[Alasan penentuan batas \emph{service}]
Batas ditarik menurut \emph{alasan berubah} dan \emph{kepemilikan data}. Sumber, Aggregator (pemilik Canonical Store),
autentikasi, penyajian, dan \emph{consumer} dipisah karena masing-masing berubah karena alasan yang berbeda. \emph{Polling},
pemetaan, penyimpanan, dan \emph{outbox} tetap berada di satu \emph{container} karena hazard dan baris \emph{outbox} harus ditulis dalam satu
transaksi. Pembagian per lapisan teknis ditolak karena mengubah transaksi itu menjadi transaksi terdistribusi, sedangkan
pembagian per sumber menjadi kandidat untuk M2. Rinciannya ada di bagian~\ref{sec:p4-batas}.
\end{jawaban}

\begin{jawaban}[Perbandingan alternatif penyimpanan]
Alternatif yang dibandingkan adalah relasional dengan kolom tetap, EAV, document \emph{store}, \emph{key-value}, dan PostgreSQL dengan
JSONB (bagian~\ref{sec:p4-storage}). PostgreSQL dengan JSONB dipilih karena memberi fleksibilitas dokumen untuk \emph{field} yang
berubah \emph{sekaligus} transaksi ACID yang dibutuhkan pola \emph{outbox}, dengan satu teknologi saja. Konsekuensinya, isi
\code{attributes} belum divalidasi tipenya dan belum diindeks.
\end{jawaban}
